\documentclass[a4paper,11pt]{article}
\usepackage[T2A]{fontenc}
\usepackage[utf8]{inputenc}
\usepackage[russian]{babel}
\usepackage{amsmath,amssymb,mathtools}
\usepackage{cmbright}
\renewcommand{\familydefault}{\sfdefault}
\usepackage{icomma}
\usepackage[a4paper,left=16mm,right=16mm,top=12mm,bottom=11mm]{geometry}
\usepackage{microtype}
\usepackage{multicol}
\usepackage{xcolor}
\usepackage{hyperref}
\definecolor{accent}{HTML}{15747A}
\definecolor{darkaccent}{HTML}{164C55}
\definecolor{textgray}{HTML}{283238}
\definecolor{muted}{HTML}{65747A}
\definecolor{rulegray}{HTML}{CADADC}
\color{textgray}
\hypersetup{hidelinks}
\setlength{\parindent}{0pt}
\setlength{\parskip}{0pt}
\setlength{\columnsep}{8mm}
\setlength{\abovedisplayskip}{1mm}
\setlength{\belowdisplayskip}{1mm}
\emergencystretch=2em
\newcommand{\sect}[1]{\par\vspace{2.2mm}{\normalsize\bfseries\color{darkaccent}#1}\par\vspace{0.3mm}{\color{rulegray}\rule{\linewidth}{0.45pt}}\par\vspace{0.8mm}}
\newcommand{\fml}[1]{\par\noindent{\normalsize\(\displaystyle #1\)}\par\vspace{2.8mm}}
\begin{document}
\thispagestyle{empty}

\hypersetup{pdftitle={Главные формулы. Ксения Клыкова},pdfauthor={Лёвин Артём Александрович}}
\begin{center}
{\small\bfseries\color{accent}УЧЕБНАЯ ПАМЯТКА \quad\textbullet\quad ОСНОВНЫЕ ФОРМУЛЫ}\par
\vspace{1mm}
{\fontsize{21}{25}\selectfont\bfseries\color{darkaccent} Вероятность, статистика и векторы\par}
\vspace{0.8mm}
{\small\color{muted}Ксения Клыкова\quad\textbullet\quad ЕГЭ · профильная математика\quad\textbullet\quad 12.10.2026}
\end{center}
\vspace{2mm}
\noindent{\color{accent}\rule{\linewidth}{0.9pt}}
\vspace{1.3mm}
{\bfseries\color{darkaccent}Главные принципы решения}\par
\vspace{1mm}
\begin{multicols}{2}
\raggedcolumns
\small
\noindent\textbullet\quad В вероятностной задаче сначала опишите события и равновозможные исходы.\par\vspace{0.75mm}
\noindent\textbullet\quad Для последовательности событий умножайте условные вероятности.\par\vspace{0.75mm}
\columnbreak
\noindent\textbullet\quad Для несовместных вариантов складывайте вероятности.\par\vspace{0.75mm}
\noindent\textbullet\quad В статистике различайте математическое ожидание, дисперсию и стандартное отклонение.\par\vspace{0.75mm}
\end{multicols}
\vspace{1.1mm}
\noindent{\color{accent}\rule{\linewidth}{0.5pt}}
\vspace{0.6mm}
\begin{multicols}{2}
\raggedcolumns
\normalsize
\sect{Правила вероятности}
\fml{P(A)=\frac{m}{n}}
\fml{P(\overline A)=1-P(A)}
\fml{P(A\cap B)=P(A)P(B\mid A)}
\fml{P(A\cup B)=P(A)+P(B)\ (A\cap B=\varnothing)}
\sect{Повторные испытания}
\fml{P_n(k)=\binom nk p^k(1-p)^{n-k}}
\fml{\binom nk=\frac{n!}{k!(n-k)!}}
\columnbreak
\sect{Статистика}
\fml{\overline x=\frac1n\sum_{i=1}^n x_i}
\fml{E(X)=\sum_i x_ip_i}
\fml{D(X)=E(X^2)-[E(X)]^2}
\fml{\sigma(X)=\sqrt{D(X)}}
\sect{Векторы}
\fml{\overrightarrow{AB}=(x_B-x_A;\,y_B-y_A)}
\fml{|\vec a|=\sqrt{a_x^2+a_y^2}}
\fml{\vec a\cdot\vec b=a_xb_x+a_yb_y}
\end{multicols}
\vfill
{\color{accent}\rule{\linewidth}{0.8pt}}\par\vspace{1.5mm}
{\small\color{muted}Лёвин Артём Александрович эксклюзивно для Ксении Клыковой}
\end{document}